
\SGARefTarget{sga3:xiv:expose}

\setcounter{footnote}{0}
\footnotetext{Version xy of 1 December 2008.}

\section{Construction of Cartan subgroups and maximal tori for a
smooth algebraic group}
\SGARefTarget{sga3:xiv:section:1}

\paragraph{\textbf{Theorem 1.1.}}
\SGARefTarget{sga3:xiv:1.1}
\emph{Let \(G\) be a smooth algebraic group over a field \(k\).
Then \(G\) admits a maximal torus \(T\), and hence a Cartan subgroup
\(C=\Centr_G(T)\).}

By \SGARefLink{sga3:xii:3.2}{XII~3.2}, finding a maximal torus \(T\) of \(G\) is equivalent to
finding a Cartan subgroup \(C\) of \(G\). Moreover, since the maximal
tori of \(G\) are those of \(G^0\), we may suppose that \(G\) is
connected. We distinguish two cases.

\medskip
\noindent\textup{1\textdegree)} \emph{The field \(k\) is finite.}
Let \(\mathcal T\) be the scheme of maximal tori of \(G\)
(\SGARefLink{sga3:xii:1.10}{XII~1.10}); it is smooth over \(k\). The group \(G\) acts on
\(\mathcal T\) through inner automorphisms, and the conjugacy theorem
(\SGARefLink{sga3:xii:6.6:item:a}{XII~6.6(a)}) says that any two \(\bar k\)-rational points of
\(\mathcal T_{\bar k}\) are conjugate under
\(G_{\bar k}(\bar k)\). Since \(\mathcal T\) is smooth over \(k\),
\(\mathcal T_{\bar k}\) is smooth over \(\bar k\); it follows that
\(\mathcal T_{\bar k}\) is isomorphic to
\[
  G_{\bar k}/N,
\]
where \(N\) is the stabilizer of an element of
\(\mathcal T_{\bar k}(\bar k)\), that is, the normalizer of a
maximal torus \(\bar T\) of \(G_{\bar k}\). Consequently
\(\mathcal T\) is a homogeneous space under the action of \(G\).
A well-known theorem of Lang
(\emph{Amer.\ J.\ Math.} \textbf{78} (1956), 555--563) says that
every homogeneous space under a smooth connected algebraic group
over a finite field \(k\) has a rational point. In particular,
\(\mathcal T\) has a rational point; equivalently, \(G\) has a
maximal torus \(T\).

\medskip
\noindent\textup{2\textdegree)} \emph{The field \(k\) is infinite.}
We shall use the following lemma.

\paragraph{\textbf{Lemma 1.2.}}
\SGARefTarget{sga3:xiv:1.2}
\emph{Let \(G\) be a smooth algebraic group over a field \(k\).
Then \(G\) admits a subgroup of type \textup{(C)}
\textup{(\SGARefLink{sga3:xiii:definition:6:2}{XIII~6.2})}.}

By \SGARefLink{sga3:xiii:corollary:6:3}{XIII~6.3}, this is equivalent to saying that \(\mathfrak g\)
contains a Cartan subalgebra \(\mathfrak d\). When \(k\) is infinite
this is immediate: \(\mathfrak g\) contains a regular element \(a\),
and one takes
\[
  \mathfrak d=\Nil(a,\mathfrak g).
\]
The case of finite \(k\) is treated exactly as in the proof of
\textup{1\textdegree)} above, but it requires first constructing the
scheme \(\mathcal D\) of Cartan subalgebras of \(\mathfrak g\) and
proving that this scheme is smooth over \(k\); this will be done
below in \SGARefLink{sga3:xiv:2.16}{2.16}.

To prove \SGARefLink{sga3:xiv:1.1}{1.1} in case \textup{2\textdegree)}, which is the case now
under consideration, it is therefore enough to know \SGARefLink{sga3:xiv:1.2}{1.2} when \(k\)
is infinite. We also record the following result for later use.

\paragraph{\textbf{Lemma 1.3.}}
\SGARefTarget{sga3:xiv:1.3}
\emph{Let \(G\) be a smooth connected affine algebraic group whose
reductive center \textup{(\SGARefLink{sga3:xii:4.1}{XII~4.1} and \SGARefLink{sga3:xiv:4.4}{4.4})} is reduced to the identity
group. Then \(G\) is nilpotent if and only if its Lie algebra
\(\mathfrak g\) is nilpotent.}

This is contained in \SGARefLink{sga3:xii:4.9}{XII~4.9}.

We can now give a procedure for constructing Cartan subgroups of
\(G\). The same procedure applies when \(k\) is finite, provided
\SGARefLink{sga3:xiv:1.2}{Lemma~1.2} is granted in that case. Suppose first that \(G\) is
affine. We argue by induction on \(n=\dim G\), the assertion being
trivial when \(n=0\). Suppose therefore that \(n>0\), and that the
assertion has been proved in dimensions \(n'<n\). Let \(Z\) be the
reductive center of \(G\), and let
\[
  u\colon G\longrightarrow G'=G/Z
\]
be the canonical homomorphism. By \SGARefLink{sga3:xii:theorem:4-7:item:c}{XII~4.7(c)}, the rule
\[
  C'\longmapsto C=u^{-1}(C')
\]
is a bijection from the Cartan subgroups of \(G'\) to those of
\(G\). Replacing \(G\) by \(G'\), we may thus suppose that the
reductive center \(Z\) of \(G\) is reduced to the identity element
\textup{(the same is true for the reductive center of \(G'\), by
\SGARefLink{sga3:xii:theorem:4-7:item:b}{XII~4.7(b)})}.

By \SGARefLink{sga3:xiv:1.2}{Lemma~1.2},\footnote{The source prints ``\SGARefLink{sga3:xiv:1.3}{1.3}'' here; the cited
existence result is \SGARefLink{sga3:xiv:1.2}{Lemma~1.2}.} \(G\) admits a subgroup \(D\) of
type \textup{(C)}. Over the algebraic closure \(\bar k\) of \(k\),
the group \(D_{\bar k}\) contains a Cartan subgroup of
\(G_{\bar k}\) \textup{(\SGARefLink{sga3:xiii:remarks:6:6:part:b}{XIII~6.6(b)})}; hence every Cartan subgroup
of \(D\) is a Cartan subgroup of \(G\) \textup{(\SGARefLink{sga3:xiii:corollary:2:8:item:a}{XIII~2.8(a)})}.
We are therefore reduced to finding a Cartan subgroup of \(D\).

If \(\dim D=\dim G\), that is, if \(D=G\), then the Lie algebra of
\(G\) is a Cartan subalgebra of itself and is therefore nilpotent.
\SGARefLink{sga3:xiv:1.3}{Lemma~1.3} then shows that \(G\) is nilpotent, so \(G\) is a Cartan
subgroup of itself. If \(\dim D<\dim G\), the induction hypothesis
gives a Cartan subgroup of \(D\), which is consequently a Cartan
subgroup of \(G\). This proves the assertion when \(G\) is affine.

In the general case, let \(Z\) be the center of \(G\). Then
\(G/Z=G'\) is affine \textup{(\SGARefLink{sga3:xii:6.1}{XII~6.1})}, and the inverse image in
\(G\) of every Cartan subgroup \(C'\) of \(G'\) is a Cartan
subgroup of \(G\) \textup{(\SGARefLink{sga3:xii:6.6:item:e}{XII~6.6(e)})}. We are thus reduced to
finding a Cartan subgroup of the smooth affine algebraic group
\(G'\), which is the case already treated.

\paragraph{\textbf{Corollary 1.4.}}
\SGARefTarget{sga3:xiv:1.4}
\emph{Let \(G\) be a smooth group scheme of finite type over an
Artinian scheme \(S\). Then \(G\) admits a maximal torus \(T\), and
hence a Cartan subgroup \(C=\Centr_G(T)\). Every torus in \(G\) is
contained in a maximal torus.}

Indeed, we may suppose that \(S\) is local, with residue field
\(k\). By \SGARefLink{sga3:xiv:1.1}{1.1}, the group
\[
  G_0=G\times_S\operatorname{Spec}(k)
\]
admits a maximal torus \(T_0\). By \SGARefLink{sga3:ix:theorem:3:6}{IX~3.6}~bis and \SGARefLink{sga3:x:corollary:2:3}{X~2.3}, \(T_0\)
comes from a torus \(T\) of \(G\), which is plainly maximal. The
last assertion follows by applying the preceding result to the
centralizer of the given torus, which is smooth over \(k\)
\textup{(\SGARefLink{sga3:xi:corollary:2:4}{XI~2.4})}.

\paragraph{\textbf{Remarks 1.5.}}
\SGARefTarget{sga3:xiv:1.5}
\textup{(a)} We shall give below, in \SGARefLink{sga3:xiv:3.20}{3.20}, \SGARefLink{sga3:xiv:3.21}{3.21}, and Exposé XV,
variants of \SGARefLink{sga3:xiv:1.4}{1.4} in which \(S\) is not assumed Artinian.

\medskip
\noindent\textup{\SGAInlineRefTarget{sga3:xiv:remark:1-5:item:b}(b)} I do not know whether every algebraic group
\(G\), not necessarily smooth, over a field \(k\) admits a maximal
torus. The question arises only in characteristic \(p>0\). Applying
\SGARefLink{sga3:xiv:1.1}{1.1} to a smooth quotient group of the form \(G'=G/I\), where \(I\)
is a suitable radicial subgroup of \(G\) \textup{(for example, the
kernel of a sufficiently high power of the Frobenius homomorphism)},
and then taking in \(G\) the inverse image of a maximal torus of
\(G'\), reduces the question to the case in which
\((G_{\bar k})_{\mathrm{red}}\) is a torus; here, as always,
\(\bar k\) denotes the algebraic closure of \(k\). The answer is
easily seen to be affirmative when \(G\) is commutative, or more
generally nilpotent: in that case \(G\) has a unique maximal torus,
which may for example be constructed by descent from the maximal
torus of \(G_{\bar k}\).\begingroup
\renewcommand{\thefootnote}{*}\footnote{M. Raynaud gave a negative
answer to the question raised here; see XVII, \SGARefLink{sga3:xvii:subsection:5-9:item:c}{Example~5.9(c)}.}
\endgroup

\medskip
\noindent\textup{(c)} When \(G\) is affine and either \(k\) is
perfect or \(G\) is solvable, \SGARefLink{sga3:xiv:1.1}{1.1} was already known, by a theorem of
Rosenlicht; his proof is quite different from the one given here.

\medskip
\noindent\textup{(d)} When \(k\) is infinite, \SGARefLink{sga3:xiv:1.1}{1.1} follows from the
much more precise result, proved below in \SGARefLink{sga3:xiv:6.1}{6.1}, that the scheme
\(\mathcal T\) of maximal tori of \(G\) is a rational variety. The
method essentially combines the proof of \SGARefLink{sga3:xiv:1.1}{1.1} with an explicit
description of the structure of the scheme \(\mathcal D\) of Cartan
subalgebras of \(\mathfrak g\). To obtain the desired result, we
must first extend to an arbitrary base prescheme certain results of
XIII~4--6 \textup{(this is the purpose of the next two sections)},
and then refine the preceding construction proving \SGARefLink{sga3:xiv:1.1}{1.1}, using the
fact that every Cartan subgroup of \(G\) is contained in one and
only one subgroup of type \textup{(C)} in \(G\) \textup{(Nos.~4
and 5)}.

\section{Lie algebras over an arbitrary prescheme: regular sections
and Cartan subalgebras}
\SGARefTarget{sga3:xiv:section:2}

\paragraph{\textbf{Lemma 2.1.}}
\SGARefTarget{sga3:xiv:2.1}
\emph{Let \(A\) be a ring, let \(\mathfrak d\) be a Lie algebra over
\(A\), and for each \(s\in\operatorname{Spec}(A)\) let
\[
  \mathfrak d(s)=\mathfrak d\otimes_A k(s)
\]
be the corresponding Lie algebra over the residue field \(k(s)\).
Suppose that the \(A\)-module \(\mathfrak d\) is of finite
presentation. The following conditions are equivalent:
\begin{enumerate}
\item[\textup{(i)}] for every \(s\in\operatorname{Spec}(A)\), the
  Lie algebra \(\mathfrak d(s)\) is nilpotent;
\item[\textup{(ii)}] for every \(x\in\mathfrak d\), the endomorphism
  \(\operatorname{ad}(x)\) is nilpotent;
\item[\textup{(iii)}] there is an integer \(N\geq 0\) such that,
  for every sequence \(x_1,\ldots,x_N\) of elements of
  \(\mathfrak d\), one has
  \[
    \operatorname{ad}(x_1)\operatorname{ad}(x_2)\cdots
    \operatorname{ad}(x_N)=0.
  \]
\end{enumerate}}
